\documentclass[12pt,a4paper]{article}
\usepackage[margin=1in]{geometry}
\usepackage[T1]{fontenc}
\usepackage[utf8]{inputenc}
\usepackage{lmodern}
\usepackage{setspace}
\usepackage{enumitem}
\usepackage{hyperref}
\usepackage{xcolor}
\hypersetup{
  colorlinks=true,
  linkcolor=blue,
  urlcolor=blue
}
\setstretch{1.15}

\title{\textbf{FitAI Trainer App}\\Complete Technical and Functional Report}
\author{Project Team}
\date{\today}

\begin{document}
\maketitle
\tableofcontents
\newpage

\section{Project Overview}
FitAI Trainer is a cross-platform fitness application built with React Native and Expo. The app provides personalized workout support, AI chat assistance, live exercise monitoring, nutrition guidance, progress tracking, and user profile management.

\subsection{Primary Goals}
\begin{itemize}[leftmargin=1.2cm]
  \item Provide an easy onboarding and authentication experience.
  \item Deliver AI-assisted guidance for fitness and nutrition.
  \item Track user workout sessions, reps, form quality, and progress.
  \item Support both online and offline usage with reliable data handling.
  \item Offer deployment for Android APK and web testing.
\end{itemize}

\section{Technology Stack}
\subsection{Frontend (Mobile + Web)}
\begin{itemize}[leftmargin=1.2cm]
  \item React Native (Expo SDK 50)
  \item React Navigation (Stack + Bottom Tabs)
  \item AsyncStorage (local persistence)
  \item Expo Camera, Expo Image Picker, Expo Speech, Expo Linear Gradient
  \item react-native-safe-area-context (safe layout on modern phones)
\end{itemize}

\subsection{Backend and AI Services}
\begin{itemize}[leftmargin=1.2cm]
  \item Python backend API (for auth, plans, stats, pose analysis, AI endpoints)
  \item Gemini-backed AI endpoints via backend proxy
  \item RESTful integration from app through centralized API client
\end{itemize}

\subsection{AI/ML Models Used (Exact)}
\begin{itemize}[leftmargin=1.2cm]
  \item \textbf{Primary LLM for chat/advice:} \texttt{gemini-2.5-flash}
  \item \textbf{Fallback LLM chain (auto):} \texttt{gemini-2.0-flash} $\rightarrow$ \texttt{gemini-2.0-flash-lite} $\rightarrow$ \texttt{gemini-1.5-flash}
  \item \textbf{Plan generation models:} \texttt{gemini-2.5-flash}, \texttt{gemini-2.0-flash}, \texttt{gemini-2.0-flash-lite}, \texttt{gemini-flash-latest} (fallback sequence)
  \item \textbf{Vision model for image analysis:} \texttt{gemini-1.5-flash} (food/exercise/general image prompts)
  \item \textbf{Pose/landmark model for live monitoring backend:} MediaPipe BlazePose (Full variant, \texttt{model\_complexity=1})
\end{itemize}

\subsection{Teacher-Friendly Explanation (Simple)}
\begin{itemize}[leftmargin=1.2cm]
  \item \textbf{Which model powers the chatbot?} The backend uses the Gemini model family. The default model is \texttt{gemini-2.5-flash}; if unavailable, the system automatically switches to the fallback chain.
  \item \textbf{How does live pose detection work?} Live monitoring does not use Gemini. It uses a computer-vision pose model (MediaPipe BlazePose). The backend extracts body keypoints (shoulders, hips, knees, ankles, etc.) from each frame and computes angles.
  \item \textbf{How are reps counted?} Each exercise has threshold angle ranges (for example, squat down/up angles). When the angle transitions through these ranges, stage transitions are detected and reps are incremented.
  \item \textbf{How are quiz-based recommendations generated?} Quiz answers (goal, level, body type, available time, health issue) are sent to the backend. The plan generator first attempts Gemini-based personalized JSON plan generation. If model generation fails, a rule-based fallback engine generates a safe plan.
  \item \textbf{Which model is used for image analysis?} Food and exercise photo analysis is processed through the backend using \texttt{gemini-1.5-flash} vision flow.
\end{itemize}

\subsection{Database and Cloud}
\begin{itemize}[leftmargin=1.2cm]
  \item Firebase Authentication (email/password flows available in project service layer)
  \item Firebase Firestore for cloud profile/session/stat data structures
  \item AsyncStorage as offline-first local storage
\end{itemize}

\section{System Architecture}
\subsection{High-Level Architecture}
\begin{enumerate}[leftmargin=1.2cm]
  \item \textbf{UI Layer:} Screens and components in React Native.
  \item \textbf{Navigation Layer:} Auth flow + Main tab flow + feature detail stacks.
  \item \textbf{Service Layer:} \texttt{src/services/api.js} and \texttt{src/services/firebase.js}.
  \item \textbf{Data Layer:} AsyncStorage local state + optional Firebase cloud sync + backend APIs.
\end{enumerate}

\subsection{Navigation Structure}
\begin{itemize}[leftmargin=1.2cm]
  \item \textbf{Startup routes:} Splash, Onboarding, Auth
  \item \textbf{Main tabs:} Dashboard, Weekly Plan, Diet, Progress, Profile
  \item \textbf{Feature stacks:} Exercise Plan, Live Exercise, Exercise Detail, Notifications, AI Chat, Admin
\end{itemize}

\section{Core Module Implementation}
\subsection{Authentication Module}
\textbf{Screen:} \texttt{AuthScreen.js}
\begin{itemize}[leftmargin=1.2cm]
  \item Supports Sign In and Sign Up modes.
  \item Uses API auth endpoints: \texttt{/auth/register}, \texttt{/auth/login}.
  \item Implements retry logic for transient network failures.
  \item On success:
  \begin{itemize}
    \item stores token in AsyncStorage,
    \item stores user profile locally,
    \item routes user to Onboarding or Main flow.
  \end{itemize}
  \item Includes offline fallback mode for delivery-safe login continuity.
\end{itemize}

\subsection{Onboarding and Plan Generation}
\textbf{Screen:} \texttt{OnboardingScreen.js}
\begin{itemize}[leftmargin=1.2cm]
  \item Captures user fitness goals, level, body type, time, and health factors.
  \item Normalizes answers into backend plan format.
  \item Calls plan endpoint: \texttt{/generate-plan}.
  \item Stores generated plan and versions in local storage for immediate app use.
\end{itemize}

\subsection{AI Chatbot Module}
\textbf{Screen:} \texttt{AIChatScreen.js}
\begin{itemize}[leftmargin=1.2cm]
  \item Multi-language conversation support (English, Urdu, Arabic).
  \item Text question handling through endpoint: \texttt{/ai-advice}.
  \item Image-based analysis through endpoint: \texttt{/analyze-image}.
  \item Model usage in backend:
  \begin{itemize}
    \item text advice/messages generated via Gemini chat model chain,
    \item image analysis requests routed to Gemini vision model.
  \end{itemize}
  \item User can send image types such as:
  \begin{itemize}
    \item Food and calorie analysis,
    \item Exercise form snapshot analysis,
    \item General fitness visual analysis.
  \end{itemize}
  \item Conversation UX includes typing states, quick suggestions, and message timeline.
\end{itemize}

\subsection{Live Exercise Monitor Module}
\textbf{Screen:} \texttt{LiveExerciseScreen.js}
\begin{itemize}[leftmargin=1.2cm]
  \item Uses device camera feed through Expo Camera.
  \item Frame analysis interval tuned by platform for performance balance.
  \item Sends image frame data to endpoint: \texttt{/analyze-pose}.
  \item Backend pose pipeline is powered by MediaPipe BlazePose for keypoint extraction and angle-based stage detection.
  \item Evaluates motion stages and angle thresholds for rep counting logic.
  \item Tracks:
  \begin{itemize}
    \item reps,
    \item exercise stage,
    \item estimated form accuracy,
    \item session timer.
  \end{itemize}
  \item Voice cues via Expo Speech for real-time coaching feedback.
  \item At session end:
  \begin{itemize}
    \item local stats updated,
    \item backend session save attempted via \texttt{/sessions},
    \item progress update event emitted to refresh dashboard/progress screens.
  \end{itemize}
\end{itemize}

\subsection{Progress and Analytics Module}
\textbf{Screen:} \texttt{ProgressScreen.js}
\begin{itemize}[leftmargin=1.2cm]
  \item Reads local and remote stats.
  \item Consumes endpoints: \texttt{/stats}, \texttt{/sessions}.
  \item Computes weekly reps and exercise distribution breakdown.
  \item Displays trend and milestone-oriented feedback.
\end{itemize}

\subsection{Diet and Nutrition Module}
\textbf{Screen:} \texttt{DietPlanScreen.js}
\begin{itemize}[leftmargin=1.2cm]
  \item Maintains food logs and nutrition calculations.
  \item Supports day-based entries and check-ins.
  \item Complements chatbot visual analysis for meal guidance.
\end{itemize}

\subsection{Profile and Preferences Module}
\textbf{Screen:} \texttt{ProfileScreen.js}
\begin{itemize}[leftmargin=1.2cm]
  \item Manages profile details and fitness preference settings.
  \item Theme preference saved in \texttt{appPrefs}.
  \item Theme updates are fully applied via controlled reload strategy.
  \item Android touch reliability improved for toggle actions.
\end{itemize}

\section{Data Management and Reliability}
\subsection{Offline-First Approach}
\begin{itemize}[leftmargin=1.2cm]
  \item Critical app data is always saved locally first via AsyncStorage.
  \item Session sync queue (\texttt{sessionSyncQueue}) stores unsent sessions if API fails.
  \item Queue flush runs automatically before session/stat reads and after successful writes.
\end{itemize}

\subsection{Request Stability}
\begin{itemize}[leftmargin=1.2cm]
  \item Unified request utility with timeout and retry support.
  \item Handles Android runtime differences where \texttt{AbortSignal.timeout} may not exist.
  \item Includes safer response parsing and structured error handling.
\end{itemize}

\subsection{Firebase Integration (Project Service Layer)}
\begin{itemize}[leftmargin=1.2cm]
  \item Firebase auth and Firestore service wrappers are implemented.
  \item Methods include user profile, onboarding answers, sessions, stats, diet logs, preferences.
  \item Hybrid helper supports local save + optional cloud sync.
\end{itemize}

\section{UI/UX and Responsiveness Implementation}
\subsection{Responsive Strategy}
\begin{itemize}[leftmargin=1.2cm]
  \item Screen-size breakpoints for small phones and tablet widths.
  \item Percentage/flexible layouts and scalable typography.
  \item Safe-area migration to \texttt{react-native-safe-area-context} across major screens.
  \item Bottom tab sizing and padding adjusted using live safe-area insets.
\end{itemize}

\subsection{Result}
\begin{itemize}[leftmargin=1.2cm]
  \item Improved header placement consistency.
  \item Better bottom control visibility on gesture-navigation devices.
  \item More stable interaction behavior on Android for profile toggles.
\end{itemize}

\section{Feature List (Delivered)}
\begin{itemize}[leftmargin=1.2cm]
  \item Authentication (Sign In / Sign Up / Guest mode)
  \item Onboarding with fitness personalization
  \item AI chatbot (text + image analysis)
  \item Live camera-based exercise monitoring with voice cues
  \item Weekly and daily workout planning
  \item Exercise detail pages with tutorials
  \item Diet logging and nutrition assistant flows
  \item Progress dashboard and statistics
  \item Notifications and guidance tips
  \item Profile management and app preferences
  \item Admin panel for controlled management flows
\end{itemize}

\section{How the App Works End-to-End}
\begin{enumerate}[leftmargin=1.2cm]
  \item User opens app, route is selected based on local token/onboarding state.
  \item User authenticates and profile context is stored.
  \item Onboarding captures user fitness parameters and generates personalized plan.
  \item User performs workouts (manual screens or live monitoring).
  \item Session events and stats are stored locally and synced to backend when available.
  \item AI chatbot assists with fitness and nutrition queries and image analysis.
  \item Progress views aggregate sessions and show metrics and milestones.
\end{enumerate}

\section{Run Instructions}
\subsection{Mobile Development (Expo)}
\begin{verbatim}
cd "d:\client working folder\hassan project\GymTrainerApp"
npm install
npx expo start
\end{verbatim}
Then scan QR with Expo Go or run emulator target.

\subsection{Android APK Build (EAS)}
\begin{verbatim}
cd "d:\client working folder\hassan project\GymTrainerApp"
npx eas login
npx eas build -p android --profile apk
\end{verbatim}

\subsection{Web/Chrome Run}
\begin{verbatim}
cd "d:\client working folder\hassan project\GymTrainerApp"
npm install
npx expo start --web
\end{verbatim}
Note: app config uses webpack for web in current setup.

\section{Backend Connectivity Notes}
\begin{itemize}[leftmargin=1.2cm]
  \item API base URL is controlled by \texttt{EXPO\_PUBLIC\_API\_URL}.
  \item Default behavior:
  \begin{itemize}
    \item Android emulator uses \texttt{10.0.2.2},
    \item iOS/web-local uses \texttt{127.0.0.1}.
  \end{itemize}
  \item For physical phone testing, API should point to LAN IP backend URL.
\end{itemize}

\section{Testing and Validation Summary}
\begin{itemize}[leftmargin=1.2cm]
  \item Navigation flow validation across auth, onboarding, tabs, and detail routes.
  \item API compile/export verification for Android bundle.
  \item Safe-area and responsive checks implemented across major screens.
  \item Profile interaction fixes validated at code level for Android touch handling.
\end{itemize}

\section{Conclusion}
The FitAI Trainer application is implemented as a production-oriented, feature-rich fitness platform with AI assistance, camera-based monitoring, personalized plan generation, and robust offline-aware data handling. The architecture supports both rapid iteration and scalable feature growth through a clean service layer, modular screens, and hybrid data persistence.

\end{document}
